\uuid{QPU0}
\titre{Approximation normale d'une loi binomiale}
\auteur{Maxime Nguyen}
\organisation{AMSCC}
\datecreate{2026-11-27}
\niveau{L2}
\module{Probabilité}
\chapitre{Probabilité continue}
\sousChapitre{Théorème central limite}
\difficulte{2}
\theme{théorème de Moivre-Laplace, loi binomiale, correction de continuité}
\serie{ordonnee}

\contenu{
\texte{
On lance $100$ fois une pièce équilibrée et on note $S$ le nombre de « pile » : $S$ suit la loi binomiale $\mathcal{B}\left(100, \frac12\right)$. On note $\mathcal{N}(\mu,\sigma)$ la loi normale d'espérance $\mu$ et d'écart-type $\sigma$, et $\Phi$ la fonction de répartition de $\mathcal{N}(0,1)$.
}

\begin{question}
\qtype{onlyone}
D'après le théorème de Moivre-Laplace, on peut approcher la loi de $S$ par la loi :
\begin{answers}
\bad{$\mathcal{N}(50, 25)$}
\good{$\mathcal{N}(50, 5)$}
\bad{$\mathcal{N}\left(100, \frac12\right)$}
\bad{$\mathcal{N}\left(50, \frac12\right)$}
\end{answers}
\begin{explanations}
On garde l'espérance $np = 50$ et l'écart-type
$\sqrt{np(1-p)} = \sqrt{25} = 5$. A donne la variance $25$ ; C et D prennent
les paramètres $n$ et $p$ de la loi binomiale pour ceux de la loi normale.

Contre-sens visé : recopier les paramètres $(n, p)$ au lieu de calculer
l'espérance et l'écart-type.
\end{explanations}
\end{question}

\begin{question}
\qtype{onlyone}
On note $Z'$ une variable de loi $\mathcal{N}(50, 5)$. Quelle approximation de $\mathbb{P}(S = 50)$ est pertinente ?
\begin{answers}
\good{$\mathbb{P}(S = 50) \approx \mathbb{P}(49{,}5 \leq Z' \leq 50{,}5) = 2\Phi(0{,}1) - 1$}
\bad{$\mathbb{P}(S = 50) \approx \mathbb{P}(Z' = 50) = 0$}
\bad{$\mathbb{P}(S = 50) = \frac12$, puisque $50$ est l'espérance de $S$}
\end{answers}
\begin{explanations}
Elle vaut environ $0{,}080$ (la valeur exacte est
environ $0{,}0796$). C'est la correction de continuité : on remplace la
valeur entière $50$ par l'intervalle $[49{,}5;50{,}5]$ de longueur $1$.

B : une variable à densité donne une probabilité nulle à chaque point ; il faut donner une épaisseur à $\{S = 50\}$.

C confond $\mathbb{P}(S = 50)$ et $\mathbb{P}(S \leq 50)$ ; même cette dernière vaut un peu plus de $\frac12$, car elle contient la valeur $50$.

Contre-sens visé : approcher un événement ponctuel d'une loi discrète par
l'événement ponctuel de la loi continue.
\end{explanations}
\end{question}

\begin{question}
\qtype{onlyone}
Avec la correction de continuité, $\mathbb{P}(S \leq 55)$ est approchée par :
\begin{answers}
\bad{$\Phi(1)$}
\bad{$\Phi(0{,}9)$}
\good{$\Phi(1{,}1)$}
\bad{$\Phi(0{,}22)$}
\end{answers}
\begin{explanations}
$\{S \leq 55\}$ contient la valeur $55$, donc on
l'approche par $\{Z' \leq 55{,}5\}$, et $\frac{55{,}5 - 50}{5} = 1{,}1$.

A est l'approximation sans correction ($55$ au lieu de $55{,}5$).

B corrige dans le mauvais sens ($54{,}5$) : cela revient à exclure la valeur $55$.

D divise par la variance $25$ au lieu de l'écart-type $5$.

Contre-sens visé : appliquer la correction de continuité dans le mauvais
sens. Le réflexe : se demander si la valeur frontière est incluse dans
l'événement.
\end{explanations}
\end{question}
}
